\chapter{Page DeepZoom : Navigation Gigapixel}
\label{chap:deepzoom}

Pour pallier les limites de la mémoire GPU lors de l'affichage de coupes 2D à ultra-haute résolution (plusieurs gigapixels), la plateforme intègre un module de visualisation tuilée (Tiled Rendering) via le composant \texttt{DeepZoomViewer} (\texttt{js/components/deepzoom-viewer.js}).
Ce système s'apparente à une navigation "Google Maps", mais appliquée à des images de microscopie massives.

\section{Le Moteur OpenSeadragon}
Le composant \texttt{DeepZoomViewer} encapsule la bibliothèque open-source \textbf{OpenSeadragon}, reconnue pour ses performances en matière de navigation fluide sur des pyramides d'images.
Le système initialise le moteur avec des paramètres optimisés pour la microscopie :
\begin{itemize}
    \item Lissage des bords de tuiles (\texttt{smoothTileEdgesMinZoom: Infinity}).
    \item Animations de zoom paramétrées.
    \item Limite stricte de chargement parallèle (\texttt{imageLoaderLimit: 8}) pour ne pas surcharger le navigateur Web.
\end{itemize}

\section{Support des Formats et Pyramides de Tuiles}
Le viewer supporte deux approches de tuilage :

\subsection{DZI (Deep Zoom Image)}
Standard classique d'OpenSeadragon. Si le jeu de données fournit un fichier \texttt{.dzi}, le viewer le délègue nativement au moteur.

\subsection{TileSource Custom (Z-Stack 2D)}
C'est le format généré de manière privilégiée par notre pipeline Python (DataPreprocessor).
Au lieu d'un format monolithique, le viewer charge un \texttt{manifest.json} (dans le dossier \texttt{tiles2d/}) qui décrit l'étendue spatiale de l'image (Width, Height, TileSize, Levels).
Ensuite, l'algorithme construit une URL dynamique en fonction du niveau de zoom (Level), du canal (Channel) et de la tranche Z (Slice) demandés par l'utilisateur :

\begin{lstlisting}[language=Java, caption=Génération d'URL dynamique pour les tuiles DeepZoom]
getTileUrl: function(level, x, y) {
    const zStr = String(sliceIndex).padStart(3, '0');
    const cStr = String(channel);
    // Ex: /DATA_WEB/.../z014_c0/level2/tile_1_3.webp
    return \texttt{${basePath}/z${zStr}\_c${cStr}/level${level}/tile\_${x}\_${y}.webp};
}
\end{lstlisting}

\section{Synchronisation et Intégration}
Bien que le DeepZoom soit techniquement un affichage 2D, il est intégré dans l'écosystème de la plateforme.
Le module expose des événements \texttt{viewport-change} et des méthodes de navigation (\texttt{setSlice}, \texttt{nudgeSlice}).
L'état de la vue (Viewport Bounds : x, y, width, height, zoom) peut être lu en temps réel via \texttt{getViewportBounds()}. Cela permet, par exemple, dans la page \textbf{Compare}, de synchroniser mathématiquement le panoramique et le zoom de deux images Gigapixel provenant de microscopes distincts, offrant ainsi une comparaison pixel-par-pixel parfaite.
